%% toc: P15 162 162.1,162.3,162.4
\chapter{Adaptedness, Sign Reversal and Stopped Information}
\label{sec:p15-w02c-p15-019}

The martingale chapter, \cref{sec:p15-martingales-and-convergence},
proves its results for one class of processes, usually the
supermartingales, and relies on two facts that are easy to leave
implicit: that each process in a martingale class is adapted, and that
the remaining classes follow by a change of sign. This chapter states
those facts as results, proves the cases that change sign, proves the
comparison step that turns event integrals into a conditional
inequality, relates the stopped field to the form used in the source,
and records which assertions of that chapter are the source's and which
are supplied in the handbook, together with the source's names and
interpretations of these objects.

\emph{Inputs.} A probability space \((\Omega,\mathcal F,P)\), a
filtration \((\mathcal F_n)_{n\geq0}\), a real process
\(X=(X_n)_{n\geq0}\), and the definitions of
\cref{def:p15-filtration-martingale,def:p15-predictable-transform,def:p15-stopping-time,def:p15-upcrossing-times}.
\emph{Conclusions.} Adaptedness is part of every martingale class and
cannot be recovered from the conditional inequality; submartingale and
martingale versions of the reduced-filtration, transform and bounded
sampling theorems, each with its own proof; a comparison lemma for
event integrals; the equality of the two forms of the stopped field;
the failure of the stopping-time property of the alternating times
for a process that is not adapted; the equivalence of the one-step and
multi-step definitions; the mean-value inequality from which the name
supermartingale comes; and the source's gambling readings, with the
doubling system on an unfavorable game and the fact that nonnegative
bounded stakes on a submartingale have nonnegative expected gain.

\section{Adaptedness belongs to the definition}
\label{sec:p15-w02c-p15-019-adapted}

\begin{hbconvention}{def:p15-w02c-p15-019-adapted-classes}{The three conditions of a martingale class}
A real process \(X\) is a martingale for \((\mathcal F_n)\) when
\begin{enumerate}
\item \(E|X_n|<\infty\) for every \(n\geq0\);
\item \(X\) is adapted: \(X_n\) is \(\mathcal F_n\)-measurable for every
      \(n\geq0\);
\item \(E(X_{n+1}\mid\mathcal F_n)=X_n\) almost surely for every \(n\geq0\).
\end{enumerate}
It is a submartingale when (iii) is replaced by
\(E(X_{n+1}\mid\mathcal F_n)\geq X_n\), and a supermartingale when it is
replaced by \(E(X_{n+1}\mid\mathcal F_n)\leq X_n\), in both cases almost
surely and with (i) and (ii) retained. These are the three conditions of
the source's definition, whose condition~(ii) is adaptedness
\cite[p.~217]{DurrettPTEAuthor2019}. Every result about the three classes
in \cref{sec:p15-martingales-and-convergence} and in the chapters that
build on it is stated for processes satisfying all three conditions.
\index{martingale!adaptedness condition}
\end{hbconvention}

Two source variants bear on this convention. The source states the
iterated inequality \(E(X_n\mid\mathcal F_m)\leq X_m\) only for
\(n>m\) \cite[thm.~4.2.4, p.~218]{DurrettPTEAuthor2019};
\cref{thm:p15-martingale-conditioning} also includes \(n=m\), and that
case is exactly condition~(ii). The source's reduced-filtration exercise
uses \(\sigma(X_1,\ldots,X_n)\) and treats martingales only
\cite[ex.~4.2.1, p.~223]{DurrettPTEAuthor2019}; the handbook includes
\(X_0\) in the natural filtration and treats all three classes
(\cref{thm:p15-w02c-p15-019-reduced-filtration}).

\begin{hblemma}{thm:p15-w02c-p15-019-multistep-definition}{proved=annex:p15-w02c-p15-019-multistep-definition}{The multi-step inequality as an equivalent definition}
Let \(X\) be adapted with \(E|X_n|<\infty\) for every \(n\). Then \(X\)
is a supermartingale if and only if
\[
        E(X_n\mid\mathcal F_m)\leq X_m\quad\text{almost surely, for all }
        n>m\geq0 .
\]
The same equivalence holds for submartingales with \(\geq\) and for
martingales with equality.
\end{hblemma}

Just before its Theorem~4.2.4 the source remarks that the conclusion of
that theorem, the inequality above for every \(n>m\), could itself be
taken as the definition of a supermartingale, and that the definition
would then be harder to check \cite[p.~218]{DurrettPTEAuthor2019}.
\Cref{thm:p15-w02c-p15-019-multistep-definition} is the equivalence
behind the remark. The difference in checking lies in what has to be
verified: the one-step form asks, for each \(n\), for one conditional
relation between \(X_n\) and the next variable, while the multi-step
form asks for one relation for every pair \(m<n\).

\begin{hbtheorem}{thm:p15-w02c-p15-019-adaptedness-needed}{proved=annex:p15-w02c-p15-019-adaptedness-needed}{What the conditional relation alone supplies}
Let \(X\) satisfy conditions (i) and (iii) of
\cref{def:p15-w02c-p15-019-adapted-classes}, without assuming (ii).
\begin{enumerate}
\item For each \(n\), \(X_n\) agrees almost surely with an integrable
      \(\mathcal F_n\)-measurable variable, namely any version of
      \(E(X_{n+1}\mid\mathcal F_n)\). If \(\mathcal F_n\) contains every
      \(P\)-null event of \(\mathcal F\), then \(X_n\) itself is
      \(\mathcal F_n\)-measurable.
\item The submartingale inequality gives no such conclusion. On
      \(\Omega=[0,1]\) with its Borel sets and Lebesgue measure, let
      \(U(\omega)=\omega\), let every \(\mathcal F_n\) be the trivial
      field \(\{\varnothing,\Omega\}\), and put
      \[
                 X_0=-2-U,\qquad X_n=-1\quad(n\geq1).
      \]
      Then \(X\) satisfies (i) and the submartingale form of (iii).
      It is not adapted, and \(X_0\) agrees almost surely with no
      \(\mathcal F_0\)-measurable variable. Moreover
      \[
           P\bigl(E(X_0\mid\mathcal F_0)<X_0\bigr)=\tfrac12 ,
      \]
      so the case \(n=m=0\) of the submartingale inequality of
      \cref{thm:p15-martingale-conditioning} fails for this process.
\end{enumerate}
\end{hbtheorem}

\section{Sign reversal and the cases it supplies}
\label{sec:p15-w02c-p15-019-sign}

\begin{hblemma}{thm:p15-w02c-p15-019-sign-duality}{proved=annex:p15-w02c-p15-019-sign-duality}{Reversing the sign of a process}
For a real process \(X\) and a filtration \((\mathcal F_n)\):
\(X\) is a submartingale if and only if \(-X\) is a supermartingale;
and \(X\) is a martingale if and only if it is both a submartingale and a
supermartingale.
The source uses both reductions in its proof of the submartingale and
martingale forms of the iterated inequality
\cite[thm.~4.2.5, p.~218]{DurrettPTEAuthor2019}.
\end{hblemma}

\begin{hblemma}{thm:p15-w02c-p15-019-event-comparison}{proved=annex:p15-w02c-p15-019-event-comparison}{From event integrals to a conditional inequality}
Let \(\mathcal G\subseteq\mathcal F\) be a sigma-algebra, let \(Y\) be
integrable, and let \(Z\) be integrable and \(\mathcal G\)-measurable.
\begin{enumerate}
\item If \(\int_AY\,dP\leq\int_AZ\,dP\) for every \(A\in\mathcal G\),
      then \(E(Y\mid\mathcal G)\leq Z\) almost surely.
\item If \(\int_AY\,dP\geq\int_AZ\,dP\) for every \(A\in\mathcal G\),
      then \(E(Y\mid\mathcal G)\geq Z\) almost surely.
\item If equality holds for every \(A\in\mathcal G\), then
      \(E(Y\mid\mathcal G)=Z\) almost surely.
\end{enumerate}
The hypothesis compares integrals over the events of \(\mathcal G\); it
does not require \(Y\leq Z\) at any point.
\end{hblemma}

\begin{hbtheorem}{thm:p15-w02c-p15-019-reduced-filtration}{proved=annex:p15-w02c-p15-019-reduced-filtration}{Reducing the filtration in each class}
Let \((\mathcal G_n)\) be a filtration with
\(\mathcal G_n\subseteq\mathcal F_n\) for every \(n\), and let \(X\) be
adapted to \((\mathcal G_n)\).
If \(X\) is a submartingale for \((\mathcal F_n)\), it is a submartingale
for \((\mathcal G_n)\). If \(X\) is a martingale for \((\mathcal F_n)\),
it is a martingale for \((\mathcal G_n)\). The supermartingale case is
the corresponding part of \cref{thm:p15-martingale-conditioning}.
For each class this applies to the natural filtration
\(\mathcal G_n=\sigma(X_0,\ldots,X_n)\).
\end{hbtheorem}

The first assertion of \cref{thm:p15-predictable-transform}, that the
transform of \(X\) by a predictable bounded \(H\) is adapted and
integrable at every finite time, places no hypothesis on \(X\). It holds
when \(X\) is adapted and \(E|X_n|<\infty\) for every \(n\), and it is to
be read with these two hypotheses. Item~1 of the next theorem proves the
assertion under them; item~2 shows that it is false when either is
dropped. Its proof in \cref{annex:p15.predictable-transform} uses both:
\(H_i\Delta X_i\) is \(\mathcal F_n\)-measurable for \(i\leq n\) because
\(H_i\) is \(\mathcal F_{i-1}\)-measurable and \(X_i\), \(X_{i-1}\) are
\(\mathcal F_i\)-measurable, and it is integrable because
\(|H_i\Delta X_i|\leq K_i(|X_i|+|X_{i-1}|)\). The second and third
assertions of that theorem assume that \(X\) is a martingale, a
submartingale or a supermartingale, and each of these classes is adapted
and integrable by conditions (i) and (ii) of
\cref{def:p15-w02c-p15-019-adapted-classes}, so they need no further
hypothesis.

\begin{hbtheorem}{thm:p15-w02c-p15-019-transform-cases}{proved=annex:p15-w02c-p15-019-transform-cases}{The transform's hypotheses and its submartingale case}
Let \(H\) be predictable with \(|H_i|\leq K_i<\infty\) for deterministic
constants \(K_i\), as in \cref{def:p15-predictable-transform}.
\begin{enumerate}
\item If \(X\) is adapted and \(E|X_n|<\infty\) for every \(n\), then
      \((H\cdot X)_n\) is \(\mathcal F_n\)-measurable and integrable for
      every finite \(n\).
\item Neither hypothesis on \(X\) in item~1 can be omitted. On
      \([0,1]\) with Lebesgue measure and \(U(\omega)=\omega\), take
      \(H_1=1\) and \(X_0=0\), with arbitrary fields and values after
      time \(1\). With \(\mathcal F_0=\mathcal F_1\) trivial
      and \(X_1=U\), the transform at time \(1\) is integrable but not
      \(\mathcal F_1\)-measurable. With \(\mathcal F_0\) trivial,
      \(\mathcal F_1=\sigma(U)\) and \(X_1=1/U\) for \(U>0\),
      \(X_1=0\) for \(U=0\), it is \(\mathcal F_1\)-measurable but not
      integrable.
\item If \(X\) is a submartingale and every \(H_i\geq0\), then
      \(H\cdot X\) is a submartingale.
\end{enumerate}
Item~3 is the case that the source's Remark after its Theorem~4.2.8
states without proof \cite[thm.~4.2.8 and remark, p.~220]{DurrettPTEAuthor2019}.
\end{hbtheorem}

\begin{hbtheorem}{thm:p15-w02c-p15-019-sampling-cases}{proved=annex:p15-w02c-p15-019-sampling-cases}{Bounded optional sampling for submartingales and martingales}
Let \(S\leq T\leq N\) be stopping times, with a deterministic integer
\(N\geq0\). If \(X\) is a submartingale, then for every
\(A\in\mathcal F_S\)
\[
               \int_AX_T\,dP\geq\int_AX_S\,dP,
               \qquad\text{hence}\qquad
               E(X_T\mid\mathcal F_S)\geq X_S
\]
almost surely. If \(X\) is a martingale, both relations are equalities.
In the source the expectation form is Theorem~4.4.1 and
Exercise~4.4.2, and the conditional form is Exercise~4.4.4, proved
there through the stopping time of Exercise~4.4.3
\cite[thm.~4.4.1, p.~235; exs.~4.4.2--4.4.4, pp.~238--239]{DurrettPTEAuthor2019};
that route is written in \cref{thm:p15-bounded-stopping-routes}.
\end{hbtheorem}

\begin{hbexample}{def:p15-w02c-p15-019-two-step}{Sampling a two-step fair walk}
Let \(\xi_1,\xi_2\) be independent fair signs, \(X_0=0\),
\(X_n=\xi_1+\cdots+\xi_n\), and
\(\mathcal F_n=\sigma(\xi_1,\ldots,\xi_n)\) with trivial
\(\mathcal F_0\). Take \(S=1\), \(N=2\), and let \(T=1\) on
\(\{\xi_1=1\}\), \(T=2\) on \(\{\xi_1=-1\}\). Then
\(\{T=1\}=\{\xi_1=1\}\in\mathcal F_1\), so \(T\) is a stopping time,
and \(S\leq T\leq N\). Here \(\mathcal F_S=\mathcal F_1\), whose events
are \(\varnothing\), \(\{\xi_1=1\}\), \(\{\xi_1=-1\}\) and \(\Omega\).

For the martingale \(X\): on \(\{\xi_1=1\}\), \(X_T=1=X_1\); on
\(\{\xi_1=-1\}\), \(X_T=-1+\xi_2\), whose integral over that event is
\(-\tfrac12\), equal to the integral of \(X_1=-1\) there. Hence
\(E(X_T\mid\mathcal F_1)=X_1\).

For the submartingale \(X^2\): on \(\{\xi_1=1\}\),
\(X_T^2=1=X_1^2\); on \(\{\xi_1=-1\}\), \(X_T^2\) is \(0\) or \(4\) with
probability \(\tfrac12\) each given \(\xi_1\), so
\(E(X_T^2\mid\mathcal F_1)=2>1=X_1^2\) there. The submartingale
inequality of \cref{thm:p15-w02c-p15-019-sampling-cases} is strict on
an event of probability \(\tfrac12\).
\end{hbexample}

\section{The stopped field and its sources}
\label{sec:p15-w02c-p15-019-stopped}

\begin{hbtheorem}{thm:p15-w02c-p15-019-stopped-field-forms}{proved=annex:p15-w02c-p15-019-stopped}{Two forms of the stopped field}
For a stopping time \(T\),
\[
 \mathcal F_T
   =\{A\in\mathcal F:A\cap\{T=n\}\in\mathcal F_n
                         \text{ for every finite }n\}.
\]
Every event \(\{T=k\}\) and \(\{T\leq k\}\), \(k\) finite, and the event
\(\{T=\infty\}\) belong to \(\mathcal F_T\), so \(T\) is
\(\mathcal F_T\)-measurable as a map into \(\{0,1,\ldots\}\cup\{\infty\}\).
The right side is the source's definition of the information known at a
stopping time \cite[p.~277]{DurrettPTEAuthor2019}, which
\cref{def:p15-stopping-time} writes with \(\{T\leq n\}\).
\end{hbtheorem}

The assertions of \cref{thm:p15-stopping-calculus} have different
sources. The equivalence of the two stopping conditions and the
preservation of the three classes under stopping are the source's
\cite[p.~220, thm.~4.2.9]{DurrettPTEAuthor2019}. That
\(\mathcal F_T\) is a sigma-algebra is not stated in the source; its
proof is \cref{annex:p15.stopping-calculus}. The inclusion
\(\mathcal F_S\subseteq\mathcal F_T\) for \(S\leq T\), and the
measurability of the stopped value, are stated in the source only in
continuous time, the latter only for Brownian motion
\cite[thms.~7.3.6 and 7.3.8, pp.~367--368]{DurrettPTEAuthor2019}. The
discrete-time statements, with their proofs, are supplied in the
handbook; the inclusion is proved by the same intersection identity as
the continuous-time theorem.

\section{Where adaptedness enters the crossing construction}
\label{sec:p15-w02c-p15-019-crossings}

The source defines the alternating times for a submartingale, so that
the process is adapted \cite[p.~220]{DurrettPTEAuthor2019}. In
\cref{def:p15-upcrossing-times} the times are defined for any real
process, and the first assertion of
\cref{thm:p15-upcrossing-inequality} holds for an adapted process: its
proof writes \(\{T_1\leq n\}\) as the union of the events
\(\{X_m\leq a\}\), \(m\leq n\), which lie in \(\mathcal F_n\) only when
\(X\) is adapted.

\begin{hbtheorem}{thm:p15-w02c-p15-019-crossing-needs-adapted}{proved=annex:p15-w02c-p15-019-crossing}{Crossing times without adaptedness}
On \([0,1]\) with Lebesgue measure and \(U(\omega)=\omega\), let every
\(\mathcal F_n\) be trivial, let \(X_n=U\) for every \(n\), and take
\(a=\tfrac12\), \(b=2\). Then \(T_1\) is not a stopping time, and
\(H_1\) is not \(\mathcal F_0\)-measurable, so \(H\) is not predictable.
\end{hbtheorem}

\section{Sources of the strategy boundary calculations}
\label{sec:p15-w02c-p15-019-sources}

In \cref{thm:p15-martingale-boundary-calculations}, item~1, the bounded
coefficients \(H_n=\xi_n\) with gain \(n\), is supplied in the handbook
as the simplest bounded strategy that is not predictable; it does not
occur in the source. Item~2 is the source's doubling system
\cite[pp.~219--220]{DurrettPTEAuthor2019}. Items~3 and~4 answer the
source's exercises asking for a martingale that tends to \(-\infty\) and
for a submartingale whose square is a supermartingale; the particular
laws are chosen in the handbook
\cite[exs.~4.2.5 and 4.2.2, p.~223]{DurrettPTEAuthor2019}.
In \cref{thm:p15-stopping-calculus,thm:p15-bounded-optional-sampling},
the locators are those given in
\cref{sec:p15-w02c-p15-019-stopped,thm:p15-w02c-p15-019-sampling-cases}.

\section{The source's names and its gambling reading}
\label{sec:p15-w02c-p15-019-readings}

This section records how the source names and interprets the objects of
\cref{sec:p15-martingales-and-convergence}. None of it is a premise of a
result there; the results stated here are proved in the annex.

\subsection*{Why a supermartingale is called super}

By \cref{thm:p15-martingale-examples}, partial sums \(S_n\) of
i.i.d.\ integrable increments with mean \(\mu\leq0\) form a
supermartingale. The source reads \(S_n\) as the fortune of a player
who bets on an unfavorable game, and observes that there is then nothing
``super'' about a supermartingale; the name comes instead from
superharmonic functions \cite[p.~217]{DurrettPTEAuthor2019}.

\begin{hbdefinition}{def:p15-w02c-p15-019-superharmonic}{Superharmonic functions and balls}
Let \(d\geq1\). A function \(f:\mathbb R^d\to\mathbb R\) is
\emph{superharmonic} if it has continuous partial derivatives of every
order at most \(2\) and
\[
   \frac{\partial^2f}{\partial x_1^2}+\cdots+\frac{\partial^2f}{\partial x_d^2}\leq0
\]
at every point. For \(x\in\mathbb R^d\) and \(r>0\),
\(B(x,r)=\{y:|x-y|\leq r\}\) is the closed ball of radius \(r\), where
\(|x-y|\) is the Euclidean distance, and \(|B(x,r)|\) is its volume,
that is, its Lebesgue measure \cite[pp.~217--218]{DurrettPTEAuthor2019}.
The source does not name the domain of \(f\); the handbook takes it to
be all of \(\mathbb R^d\), so that every ball lies in it.
\index{superharmonic function}
\end{hbdefinition}

\begin{hbtheorem}{thm:p15-w02c-p15-019-mean-value}{proved=annex:p15-w02c-p15-019-mean-value}{The mean-value inequality for superharmonic functions}
If \(f\) is superharmonic on \(\mathbb R^d\), then for every
\(x\in\mathbb R^d\) and every \(r>0\)
\[
        f(x)\geq\frac1{|B(x,r)|}\int_{B(x,r)}f(y)\,dy .
\]
This is the source's display~(4.2.1), which it states without proof
\cite[p.~217, eq.~(4.2.1)]{DurrettPTEAuthor2019}.
\end{hbtheorem}

The supermartingale inequality \(E(X_{n+1}\mid\mathcal F_n)\leq X_n\)
has the same shape as (4.2.1): the present value is at least an average
of the values that may follow it, with the conditional mean in the place
of the average over a ball. This comparison is the handbook's reading of
the source's remark; the source states only that the name comes from
the inequality.

\begin{hbexample}{def:p15-w02c-p15-019-mean-value-example}{A strict case and an equality case}
Take \(d=1\), so \(B(x,r)=[x-r,x+r]\) and \(|B(x,r)|=2r\). The function
\(f(y)=-y^2\) has \(f''=-2\leq0\), and
\[
 \frac1{2r}\int_{x-r}^{x+r}(-y^2)\,dy
   =-\frac{(x+r)^3-(x-r)^3}{6r}
   =-x^2-\frac{r^2}3<f(x),
\]
so the inequality of \cref{thm:p15-w02c-p15-019-mean-value} is strict
at every \(x\) and \(r\). For \(f(y)=cy\) with a real constant \(c\) the
second derivative is \(0\), both \(f\) and \(-f\) are superharmonic, and
the average over \([x-r,x+r]\) is exactly \(cx=f(x)\): equality, the
analogue of a martingale.
\end{hbexample}

\subsection*{Stakes, winnings and shares}

The source reads a predictable sequence as a gambling system
\cite[p.~219]{DurrettPTEAuthor2019}. \(H_n\) is the amount a gambler
bets at time \(n\); it may use the outcomes at times
\(1,\ldots,n-1\) but not the outcome at time \(n\), which is what
\(\mathcal F_{n-1}\)-measurability expresses. If \(X_n\) is the net
amount won by time \(n\) by a player who bets one dollar every time,
then \((H\cdot X)_n\) is the net amount won with the system \(H\): a
wager of three dollars at time \(m\) changes the fortune by three times
the change for a wager of one dollar, so the change at time \(m\) is
\(H_m(X_m-X_{m-1})\). The source gives a second reading in which
\(X_m\) is the value of a stock and \(H_m\) is the number of shares held
from time \(m-1\) to time \(m\); the term \(H_m(X_m-X_{m-1})\) is then
the change in the value of the holding over that period.

\subsection*{The system called the martingale}

The doubling system of item~2 of
\cref{thm:p15-martingale-boundary-calculations} (stake one dollar first,
double the stake after each loss, return to one dollar after a win) is
the famous gambling system that the source calls ``the martingale''
\cite[pp.~219--220]{DurrettPTEAuthor2019}. A run of \(k\) losses
followed by a win ends with a net gain of one dollar; the source's table
of four losses and a win is reproduced in
\cref{annex:p15.martingale-boundary-calculations}. For this reason the
system seems to provide a sure thing as long as \(P(\xi_m=1)>0\). The
source answers with its Theorem~4.2.8, the transform theorem
\cref{thm:p15-predictable-transform}: there is no system for beating an
unfavorable game \cite[p.~220]{DurrettPTEAuthor2019}. The next result
makes the answer exact for this system.

\begin{hbcorollary}{thm:p15-w02c-p15-019-doubling-unfavorable}{proved=annex:p15-w02c-p15-019-doubling-unfavorable}{The doubling system on an unfavorable game}
Let \(\xi_1,\xi_2,\ldots\) be independent with
\(P(\xi_n=1)=p\) and \(P(\xi_n=-1)=1-p\), where \(0<p\leq\tfrac12\);
let \(X_0=0\), \(X_n=\xi_1+\cdots+\xi_n\),
\(\mathcal F_0=\{\varnothing,\Omega\}\) and
\(\mathcal F_n=\sigma(\xi_1,\ldots,\xi_n)\). Let \(H\) be the doubling
system, \(T=\inf\{n\geq1:\xi_n=1\}\) and \(G_n=(H\cdot X)_{T\wedge n}\).
Then \(H\) is predictable with \(1\leq H_n\leq2^{n-1}\); \(H\cdot X\) and
\(G\) are supermartingales starting at \(0\), so
\[
          E(H\cdot X)_n\leq0\quad\text{and}\quad EG_n\leq0
          \qquad\text{for every }n;
\]
and yet \(G_n\to1\) almost surely. For \(p=\tfrac12\) both processes are
martingales and both expectations are \(0\).
\end{hbcorollary}

The exact value \(EG_n=1-[2(1-p)]^n\) of
\cref{thm:p15-martingale-boundary-calculations} agrees: it is at most
\(0\) exactly when \(p\leq\tfrac12\). The sure gain is collected only at
the random time \(T\), which has no deterministic bound, and on the
event \(\{T>n\}\) the loss \(2^n-1\) grows without bound in \(n\).

The source draws a second reading from the same theorem: although it
implies that one cannot make money with gambling systems, gambling
systems can be used to prove theorems \cite[p.~220]{DurrettPTEAuthor2019}.
Its first instance is the upcrossing system of
\cref{def:p15-upcrossing-times}, which proves
\cref{thm:p15-upcrossing-inequality}.

\subsection*{Stopping as a gambling system}

The source introduces stopping times through a very special system:
bet one dollar at each time \(n\leq N\), then stop playing
\cite[p.~220]{DurrettPTEAuthor2019}. In the notation of
\cref{def:p15-stopping-time}, with \(T\) for the source's \(N\), the
stakes are \(H_n=\mathbf1_{\{T\geq n\}}\). They are predictable because
\(\{T\geq n\}=\{T\leq n-1\}^c\in\mathcal F_{n-1}\), and the winnings are
\((H\cdot X)_n=X_{T\wedge n}-X_0\), the identity displayed in
\cref{thm:p15-stopping-calculus}. With the transform theorem this is the
source's proof that a stopped supermartingale is a supermartingale
\cite[thm.~4.2.9, p.~220]{DurrettPTEAuthor2019}.

\subsection*{The key fact behind the upcrossing inequality}

After proving the upcrossing inequality, the source remarks that its
classical form is a little misleading. The key fact in the proof is
\(E(K\cdot Y)_n\geq0\) for the clipped submartingale \(Y\) and the
complementary stakes \(K_m=1-H_m\): however hard one tries, one cannot
lose money betting on a submartingale
\cite[p.~221]{DurrettPTEAuthor2019}. The next result states that fact
for every system of bounded nonnegative stakes.

\begin{hbcorollary}{thm:p15-w02c-p15-019-no-loss}{proved=annex:p15-w02c-p15-019-no-loss}{Betting on a submartingale}
Let \(Y\) be a submartingale and let \(K\) be predictable with
\(0\leq K_m\leq c_m\) for deterministic constants \(c_m<\infty\). Then
\(K\cdot Y\) is a submartingale with \((K\cdot Y)_0=0\), and
\(E(K\cdot Y)_n\geq0\) for every \(n\).
\end{hbcorollary}

The proof of \cref{thm:p15-upcrossing-inequality} uses this with
\(K_m=1-H_m\), where \(H\) is the upcrossing system; then
\((H\cdot Y)_n=Y_n-Y_0-(K\cdot Y)_n\) gives
\(E(H\cdot Y)_n\leq E(Y_n-Y_0)\), and the bound
\((b-a)U_n(a,b)\leq(H\cdot Y)_n\) gives the inequality.

\section{Relationships}
\label{sec:p15-w02c-p15-019-relationships}

\Cref{thm:p15-w02c-p15-019-sign-duality} turns every supermartingale
proof in \cref{sec:p15-martingales-and-convergence} into a
submartingale proof applied to \(-X\); the reduced-filtration, transform
and sampling theorems above write those cases out, the last through
\cref{thm:p15-w02c-p15-019-event-comparison}. The same comparison lemma
completes the conditional step of
\cref{thm:p15-bounded-optional-sampling,thm:p15-bounded-stopping-routes}.
The convention \cref{def:p15-w02c-p15-019-adapted-classes} is the one
stated at the start of the optional-stopping chapter in
\cref{def:p15-os-conventions}.
\Cref{thm:p15-w02c-p15-019-adaptedness-needed} shows why it must be
stated for sub- and supermartingales: for a martingale the identity
already forces adaptedness up to null events, but an inequality does
not.
\Cref{thm:p15-w02c-p15-019-no-loss} is the submartingale case of
\cref{thm:p15-w02c-p15-019-transform-cases} combined with mean
monotonicity, and \cref{thm:p15-w02c-p15-019-doubling-unfavorable}
combines the supermartingale case of \cref{thm:p15-predictable-transform}
with the exact gains of \cref{thm:p15-martingale-boundary-calculations}.
\Cref{thm:p15-w02c-p15-019-mean-value} is used only to explain a name;
no probabilistic result depends on it.
The conditional-limit theorems of
\cref{sec:p15-uniform-integrability-and-conditional-limits} use
condition~(ii) of \cref{def:p15-w02c-p15-019-adapted-classes} in three
places: it makes the limit of a uniformly integrable submartingale
measurable with respect to \(\sigma(\bigcup_n\mathcal F_n)\), it puts in
\(\mathcal F_n\) the events on which \(X_n\) exceeds
\(E(X_\infty\mid\mathcal F_n)\) by at least \(1/k\), and it makes \(X_n\)
itself a version of \(E(X_\infty\mid\mathcal F_n)\) when \(X_n\to X_\infty\)
in \(L^1\) for a martingale
(\cref{thm:p15-submartingale-ui,thm:p15-closed-martingale}).
